\expandafter\gdef\csname SNside1\endcsname{\node[anchor=north west,inner sep=0pt,text width=48mm,font=\SideFont,align=left] at ([xshift=138mm,yshift=-500.62pt]current page.north west) {\textcolor{SNink}{\textbf{最小反例 · 相同数字隐藏记录身份}}\par\vspace{1.5mm}最小反例给相等关键字附身份：\(2_a,2_b,1\)经简单选择首趟成为\(1,2_b,2_a\)。只写数字2、2会看不出顺序改变；稳定性讨论的是整条记录的身份。\par\smallskip\begin{tikzpicture}[x=1mm,y=1mm,>=Stealth,node distance=5mm]
\node[draw=SNink!50,fill=SNblue,rounded corners=1mm,text width=34mm,align=center,inner sep=1.2mm,font=\fontsize{7.5}{9}\selectfont] (m0) at (0,0) {2a};
\node[draw=SNink!50,fill=SNblue,rounded corners=1mm,text width=34mm,align=center,inner sep=1.2mm,font=\fontsize{7.5}{9}\selectfont] (m1) at (0,-11) {2b};
\node[draw=SNink!50,fill=SNblue,rounded corners=1mm,text width=34mm,align=center,inner sep=1.2mm,font=\fontsize{7.5}{9}\selectfont] (m2) at (0,-22) {1};
\draw[->,SNink] (m0) edge[bend left=40] (m2);
\end{tikzpicture}
\par{\fontsize{7.5}{9}\selectfont 交换首项与1跨过2b，结果为1、2b、2a。}};}
\expandafter\gdef\csname SNside2\endcsname{\node[anchor=north west,inner sep=0pt,text width=48mm,font=\SideFont,align=left] at ([xshift=138mm,yshift=-348.56pt]current page.north west) {\textcolor{SNink}{\textbf{不变量 · 相等时继续向右定位}}\par\vspace{1.5mm}折半插入寻找首个大于x的位置，旧相等项都留在左边。定位快只减少比较，数组腾位置仍要移动；看到折半二字就写整体对数时间，漏掉了主要移动量。};}
\expandafter\gdef\csname SNside3\endcsname{\node[anchor=north west,inner sep=0pt,text width=48mm,font=\SideFont,align=left] at ([xshift=138mm,yshift=-95.00pt]current page.north west) {\textcolor{SNink}{\textbf{错项排除 · 最好线性依赖早停标志}}\par\vspace{1.5mm}起泡有序输入一趟无交换即可停，所以线性；若程序固定执行全部趟数，即使每次都不交换，比较仍是二次。复杂度表需要与实际程序条件配套。};}
\expandafter\gdef\csname SNside4\endcsname{\node[anchor=north west,inner sep=0pt,text width=48mm,font=\SideFont,align=left] at ([xshift=138mm,yshift=-95.00pt]current page.north west) {\textcolor{SNink}{\textbf{不变量 · d间隔有序不是相邻有序}}\par\vspace{1.5mm}一趟希尔处理全部同余类，得到每个间隔d子序列有序。不同余数类之间仍未比较；只有最后d=1，才把全数组约束为同一有序序列。};}
\expandafter\gdef\csname SNside5\endcsname{\node[anchor=north west,inner sep=0pt,text width=48mm,font=\SideFont,align=left] at ([xshift=138mm,yshift=-95.00pt]current page.north west) {\textcolor{SNink}{\textbf{错项排除 · 首枢轴遇到已有序可能最坏}}\par\vspace{1.5mm}首记录作枢轴时，升序输入每次都产生0与n\ensuremath{-}1的分割，时间二次、栈线性。插入排序有序最快的经验不能迁移到这份快排；随机枢轴需另用期望模型分析。};}
\expandafter\gdef\csname SNside6\endcsname{\node[anchor=north west,inner sep=0pt,text width=48mm,font=\SideFont,align=left] at ([xshift=138mm,yshift=-95.00pt]current page.north west) {\textcolor{SNink}{\textbf{不变量 · 尾部已序，前部只是堆}}\par\vspace{1.5mm}升序堆排序每轮把当前最大值固定到尾部，下一轮只调整缩小后的前部堆。前部兄弟不保证有序，不能把一次建堆的数组误判为一趟完整升序输出。};}
\expandafter\gdef\csname SNside7\endcsname{\node[anchor=north west,inner sep=0pt,text width=48mm,font=\SideFont,align=left] at ([xshift=138mm,yshift=-95.00pt]current page.north west) {\textcolor{SNink}{\textbf{课程连接 · 优先队列与选择小部分记录}}\par\vspace{1.5mm}只找最小10项时可维护10项大根堆，把更小新值替换堆顶；不必完整排序。流式读入适合这种状态维护，若可原地随机访问，还可比较选择算法的期望线性方案。};}
\expandafter\gdef\csname SNside8\endcsname{\node[anchor=north west,inner sep=0pt,text width=48mm,font=\SideFont,align=left] at ([xshift=138mm,yshift=-95.00pt]current page.north west) {\textcolor{SNink}{\textbf{错项排除 · 比较上界不等于具体次数}}\par\vspace{1.5mm}归并a、b项最坏比较a+b\ensuremath{-}1次，但一段提前耗尽就直接复制另一段。计算给定值的过程题应逐次比较段首，不能每次都机械代入最坏上界。};}
\expandafter\gdef\csname SNside9\endcsname{\node[anchor=north west,inner sep=0pt,text width=48mm,font=\SideFont,align=left] at ([xshift=138mm,yshift=-95.00pt]current page.north west) {\textcolor{SNink}{\textbf{不变量 · LSD依赖每轮的稳定性}}\par\vspace{1.5mm}更高位相同的记录仍要靠已经排好的低位决定先后，故同桶必须保留上轮次序。改成插桶头或逆序收集可能破坏这个不变式，即使每轮单独看都按本位有序。};}
\expandafter\gdef\csname SNside10\endcsname{\node[anchor=north west,inner sep=0pt,text width=48mm,font=\SideFont,align=left] at ([xshift=138mm,yshift=-95.00pt]current page.north west) {\textcolor{SNink}{\textbf{课程连接 · 外排瓶颈联系操作系统I/O}}\par\vspace{1.5mm}外排优先按块读写，输出缓冲攒满再写可减少细粒度I/O。多路需要更多输入缓冲，路数受实际内存约束；双缓冲与异步I/O还会改变可用块数和耗时模型。};}
\expandafter\gdef\csname SNside11\endcsname{\node[anchor=north west,inner sep=0pt,text width=48mm,font=\SideFont,align=left] at ([xshift=138mm,yshift=-95.00pt]current page.north west) {\textcolor{SNink}{\textbf{最小反例 · 输出段可以比内存更长}}\par\vspace{1.5mm}最小直观情形：工作区1条，输入1、2、3。每输出一条再换入更大者，三条属于同一段，内存始终只存1条。不能把段长固定等同工作区容量。};}
\expandafter\gdef\csname SNside12\endcsname{\node[anchor=north west,inner sep=0pt,text width=48mm,font=\SideFont,align=left] at ([xshift=138mm,yshift=-95.00pt]current page.north west) {\textcolor{SNink}{\textbf{错项排除 · 下界没有否定计数排序}}\par\vspace{1.5mm}比较排序下界假设只凭两两比较区分互异关键字排列；计数索引利用整数取值范围，成本还含k。若k很大，线性于n的口号忽略了计数数组和前缀和代价。};}
\expandafter\gdef\csname SNside13\endcsname{\node[anchor=north west,inner sep=0pt,text width=48mm,font=\SideFont,align=left] at ([xshift=138mm,yshift=-95.00pt]current page.north west) {\textcolor{SNink}{\textbf{最小反例 · 输出段可以比内存更长}}\par\vspace{1.5mm}最小直观情形：工作区1条，输入1、2、3。每输出一条再换入更大者，三条属于同一段，内存始终只存1条。不能把段长固定等同工作区容量。};}
\expandafter\gdef\csname SNside14\endcsname{\node[anchor=north west,inner sep=0pt,text width=48mm,font=\SideFont,align=left] at ([xshift=138mm,yshift=-95.00pt]current page.north west) {\textcolor{SNink}{\textbf{错项排除 · 下界没有否定计数排序}}\par\vspace{1.5mm}比较排序下界假设只凭两两比较区分互异关键字排列；计数索引利用整数取值范围，成本还含k。若k很大，线性于n的口号忽略了计数数组和前缀和代价。};}
\expandafter\gdef\csname SNside15\endcsname{\node[anchor=north west,inner sep=0pt,text width=48mm,font=\SideFont,align=left] at ([xshift=138mm,yshift=-495.93pt]current page.north west) {\textcolor{SNink}{\textbf{不变量 · LSD依赖每轮的稳定性}}\par\vspace{1.5mm}更高位相同的记录仍要靠已经排好的低位决定先后，故同桶必须保留上轮次序。改成插桶头或逆序收集可能破坏这个不变式，即使每轮单独看都按本位有序。};}
\expandafter\gdef\csname SNside16\endcsname{\node[anchor=north west,inner sep=0pt,text width=48mm,font=\SideFont,align=left] at ([xshift=138mm,yshift=-95.00pt]current page.north west) {\textcolor{SNink}{\textbf{课程连接 · 外排瓶颈联系操作系统I/O}}\par\vspace{1.5mm}外排优先按块读写，输出缓冲攒满再写可减少细粒度I/O。多路需要更多输入缓冲，路数受实际内存约束；双缓冲与异步I/O还会改变可用块数和耗时模型。};}
\AddToShipoutPictureFG{\begin{tikzpicture}[remember picture,overlay]\ifcsname SNside\thepage\endcsname\csname SNside\thepage\endcsname\fi\end{tikzpicture}}